\chapter{Cohomology and base change}
\label{mk-ch-cohomology}

For a \(k\)-scheme \(T\), put \(C_T=C\times_kT\).  Cohomology is taken on
the Zariski site (equivalently for quasi-coherent sheaves on the small fppf
site).

\begin{theorem}[Cohomology and base change]
  \label{mk-thm-cohomology-base-change}
  \dcref{III.12,custom}
  \uses{mk-def-genus}
  \cite{stacks-project,stacks-cohomology-schemes}
  For a coherent sheaf \(\mathcal F\) on \(C\), \(H^i(C,\mathcal F)\) is
  finite-dimensional over \(k\).  If \(k\to k'\) is a field extension, then
  \[
    H^i(C,\mathcal F)\otimes_k k'\cong
    H^i(C_{k'},\mathcal F_{k'}).
  \]
\end{theorem}

\begin{proposition}[Global functions]
  \label{mk-prop-global-functions}
  \dcref{III.1,custom}
  \cite{hartshorne-ag,stacks-schemes}
  If \(C\) is geometrically integral and proper over \(k\), then
  \(H^0(C,\mathcal O_C)=k\).
\end{proposition}

\begin{proposition}[The structure sheaf controls the tangent dimension]
  \label{mk-prop-genus-tangent}
  \dcref{5.11,custom}
  \uses{mk-def-genus,mk-thm-cohomology-base-change}
  For the identity of the Picard scheme, the Zariski tangent space satisfies
  \[
    T_0\operatorname{Pic}_{C/k}\cong H^1(C,\mathcal O_C).
  \]
  In particular its dimension is \(g(C)\), and this identification commutes
  with extension of the ground field.
\end{proposition}
